%% Guided notes, AP Precalculus Unit 1 (lesson plan: 1.1-1.6, 1.7, 1.8, 1.9-1.10, 1.11, 1.12, 1.13-1.14).
%% \fillin{...} = blank for students / answer in the filled version; workspace = empty box / worked solution.

%%DOC 1.1 | Change in Tandem
\goals{\item Describe how output values change as input values increase.
\item Identify intervals where a function is increasing or decreasing, and where its graph is concave up or concave down.
\item Identify zeros and the sign of a function from a graph, a table or a formula.}
\vocab{function, input, output, increasing, decreasing, rate of change, concave up, concave down, zero, $x$-intercept, positive, negative.}

\subsection*{1. What is a function?}
A \textbf{function} is a rule that assigns to each \fillin{input value} exactly \fillin{one} output value. We write $y=f(x)$: $x$ is the \fillin{input (independent variable)} and $f(x)$ is the \fillin{output (dependent variable)}. As the input values change, the output values change \emph{in tandem}.

\subsection*{2. Increasing and decreasing}
\begin{keyidea}
On an interval, $f$ is \textbf{increasing} if, as the input values increase, the output values always \fillin{increase}: if $a<b$ then $f(a)$ \fillin{$<$} $f(b)$.\\
$f$ is \textbf{decreasing} if, as the input values increase, the output values always \fillin{decrease}: if $a<b$ then $f(a)$ \fillin{$>$} $f(b)$.\\
Intervals of increase and decrease are written with \fillin{input ($x$)} values.
\end{keyidea}

\subsection*{3. Concavity: how is the function changing?}
\begin{keyidea}
The graph of $f$ is \textbf{concave up} on an interval if the rate of change of $f$ is \fillin{increasing}: the graph bends \fillin{upward}.\\
The graph is \textbf{concave down} if the rate of change is \fillin{decreasing}: the graph bends \fillin{downward}.
\end{keyidea}
\begin{center}
\begin{tikzpicture}
\begin{axis}[apaxis,xmin=-0.3,xmax=4.3,ymin=-0.3,ymax=4.5,width=4.4cm,height=4cm,xtick=\empty,ytick=\empty,title={\scriptsize increasing, concave up}]
\addplot[apblue,very thick,domain=0:4]{0.25*x^2};\end{axis}\end{tikzpicture}
\begin{tikzpicture}
\begin{axis}[apaxis,xmin=-0.3,xmax=4.3,ymin=-0.3,ymax=4.5,width=4.4cm,height=4cm,xtick=\empty,ytick=\empty,title={\scriptsize increasing, concave down}]
\addplot[apblue,very thick,domain=0:4]{2*sqrt(x)};\end{axis}\end{tikzpicture}
\begin{tikzpicture}
\begin{axis}[apaxis,xmin=-0.3,xmax=4.3,ymin=-0.3,ymax=4.5,width=4.4cm,height=4cm,xtick=\empty,ytick=\empty,title={\scriptsize decreasing, concave up}]
\addplot[apblue,very thick,domain=0:4]{4/(x+1)};\end{axis}\end{tikzpicture}
\begin{tikzpicture}
\begin{axis}[apaxis,xmin=-0.3,xmax=4.3,ymin=-0.3,ymax=4.5,width=4.4cm,height=4cm,xtick=\empty,ytick=\empty,title={\scriptsize decreasing, concave down}]
\addplot[apblue,very thick,domain=0:4]{4-0.25*x^2};\end{axis}\end{tikzpicture}
\end{center}

\subsection*{4. Zeros and sign}
A \textbf{zero} of $f$ is an input $a$ such that \fillin{$f(a)=0$}. Graphically, the real zeros are the \fillin{$x$-intercepts}. The function is \textbf{positive} where its graph is \fillin{above} the $x$-axis and \textbf{negative} where it is \fillin{below}. A function whose graph has no breaks can change sign only at a \fillin{zero}.

\ex{Example 1 \textbullet\ Reading a graph}
\begin{center}
\begin{tikzpicture}
\begin{axis}[gridaxis,xmin=-2.5,xmax=2.5,ymin=-3.5,ymax=3.5,width=6.5cm,height=5.5cm]
\addplot[apblue,very thick,domain=-2.2:2.2,samples=120]{-x^3+3*x};
\node[circle,fill=apblue,inner sep=1.5pt] at (axis cs:-1,-2){};\node[circle,fill=apblue,inner sep=1.5pt] at (axis cs:1,2){};
\end{axis}
\end{tikzpicture}
\end{center}
The graph of $f(x)=-x^3+3x$ is shown. Find (a) the intervals where $f$ is increasing and decreasing, (b) where the graph is concave up and concave down, (c) the zeros, (d) where $f(x)>0$.
\begin{workspace}[4.2cm]
(a) Increasing on $(-1,1)$; decreasing on $(-\infty,-1)$ and $(1,\infty)$.\\
(b) Concave up on $(-\infty,0)$; concave down on $(0,\infty)$. The change happens at the point $(0,0)$.\\
(c) $-x^3+3x=-x(x^2-3)=0$ gives $x=0$ and $x=\pm\sqrt3\approx\pm1.73$.\\
(d) $f(x)>0$ on $(-\infty,-\sqrt3)\cup(0,\sqrt3)$.
\end{workspace}

\ex{Example 2 \textbullet\ Reading a table}
\begin{center}\begin{tabular}{c|cccccc}$x$&$0$&$1$&$2$&$3$&$4$&$5$\\\hline$g(x)$&$-6$&$-1$&$2$&$3$&$2$&$-1$\end{tabular}\end{center}
Describe how $g$ changes. Where must $g$ have zeros (assume the graph of $g$ has no breaks)?
\begin{workspace}[3.6cm]
The outputs increase from $x=0$ to $x=3$, then decrease. The changes $5,3,1,-1,-3$ decrease steadily, so the graph is concave down.\\
The sign changes between $x=1$ and $x=2$ and between $x=4$ and $x=5$: $g$ has a zero in each of these intervals.
\end{workspace}

\ex{Example 3 \textbullet\ In context}
The temperature $T$ ($^\circ$C) of a cup of coffee is $90, 70, 56, 46, 40$ at $t=0,5,10,15,20$ minutes. Describe $T$.
\begin{workspace}[2.8cm]
$T$ is decreasing (the coffee cools). The changes are $-20,-14,-10,-6$: the temperature falls more and more slowly, so the rate of change is increasing and the graph is concave up.
\end{workspace}

\begin{mistakes}
\begin{itemize}
\item Writing intervals with $y$-values instead of $x$-values.
\item Confusing ``decreasing'' with ``concave down''. A graph can be increasing and concave down at the same time.
\end{itemize}
\end{mistakes}

\begin{quickcheck}
1. Outputs at $x=0,1,2,3$ are $10,7,5,4$. The function is \fillin{decreasing} and its graph is concave \fillin{up}.\\
2. The zeros of $f(x)=(x-2)(x+5)$ are \fillin{$x=2$ and $x=-5$}.
\end{quickcheck}

%%DOC 1.2 | Rates of Change
\goals{\item Compute the average rate of change of a function over an interval, from a formula, a table or a graph.
\item Interpret an average rate of change in context, with units.
\item Estimate the rate of change at a point, and relate its sign to increasing and decreasing behavior.}
\vocab{average rate of change, secant line, slope, rate of change at a point, units.}

\subsection*{1. Average rate of change}
\begin{keyidea}
The average rate of change of $f$ over the interval $[a,b]$ is
\[\frac{\fillin{$f(b)-f(a)$}}{\fillin{$b-a$}}.\]
It is the slope of the \fillin{secant} line through the points $\big(a,f(a)\big)$ and $\big(b,f(b)\big)$. Its units are \fillin{output units} per \fillin{input unit}.
\end{keyidea}

\begin{center}
\begin{tikzpicture}
\begin{axis}[gridaxis,xmin=-0.5,xmax=4.5,ymin=-0.5,ymax=9.5,ytick distance=2,width=6cm,height=5.5cm]
\addplot[apblue,very thick,domain=0:3.2]{x^2};
\addplot[apred,thick,domain=0.3:3.6]{4*x-3};
\node[circle,fill=apred,inner sep=1.6pt] at (axis cs:1,1){};\node[circle,fill=apred,inner sep=1.6pt] at (axis cs:3,9){};
\node[font=\small,apred] at (axis cs:3.9,7.8){secant};
\end{axis}
\end{tikzpicture}
\end{center}

\ex{Example 1 \textbullet\ From a formula}
Find the average rate of change of $f(x)=x^2$ on $[1,3]$, and of $g(x)=x^3-4x$ on $[-2,1]$.
\begin{workspace}[3cm]
$\dfrac{f(3)-f(1)}{3-1}=\dfrac{9-1}{2}=4$.\qquad $g(1)=-3$, $g(-2)=-8+8=0$: $\dfrac{-3-0}{1-(-2)}=-1$.
\end{workspace}

\ex{Example 2 \textbullet\ From a graph}
\begin{center}
\begin{tikzpicture}
\begin{axis}[gridaxis,xmin=-0.5,xmax=6.5,ymin=-0.5,ymax=5.5,width=6cm,height=4.8cm]
\addplot[apblue,very thick] coordinates {(0,1)(2,5)(4,4)(6,0)};
\end{axis}
\end{tikzpicture}
\end{center}
Find the average rate of change of $h$ on $[0,2]$, $[2,6]$ and $[0,6]$.
\begin{workspace}[2.6cm]
$[0,2]$: $\frac{5-1}{2}=2$. \quad $[2,6]$: $\frac{0-5}{4}=-\frac54$. \quad $[0,6]$: $\frac{0-1}{6}=-\frac16$.
\end{workspace}

\subsection*{2. Interpreting in context}
\begin{keyidea}[Template]
``On average, [quantity] \fillin{increases or decreases} by [number] [output units] per [input unit] from [input $a$] to [input $b$].''
\end{keyidea}

\ex{Example 3 \textbullet\ Units}
$V(t)$ is the number of visitors to a website $t$ hours after midnight; $V(2)=150$ and $V(6)=870$. Compute and interpret $\frac{V(6)-V(2)}{6-2}$.
\begin{workspace}[2.4cm]
$\frac{870-150}{4}=180$. From 2 a.m.\ to 6 a.m., the number of visitors increased on average by 180 visitors per hour.
\end{workspace}

\subsection*{3. Rate of change at a point}
The rate of change at $x=a$ can be \fillin{estimated} by average rates of change over \fillin{smaller and smaller} intervals that contain $a$, using values on both sides.

\ex{Example 4 \textbullet\ Estimating}
Estimate the rate of change of $f(x)=x^3$ at $x=2$.
\begin{workspace}[3cm]
$[2,2.1]$: $\frac{9.261-8}{0.1}=12.61$. \ $[2,2.01]$: $12.0601$. \ $[1.99,2]$: $\frac{8-7.880599}{0.01}=11.9401$.\\
The values approach $12$: the rate of change at $x=2$ is about $12$.
\end{workspace}

\begin{keyidea}[Sign of the rate of change]
If the rate of change is positive on an interval, $f$ is \fillin{increasing} there; if it is negative, $f$ is \fillin{decreasing}. At a point where $f$ changes from increasing to decreasing, the rate of change changes from \fillin{positive to negative}.
\end{keyidea}

\begin{mistakes}
\begin{itemize}
\item Subtracting in different orders: $\frac{f(b)-f(a)}{a-b}$.
\item Giving the change $f(b)-f(a)$ instead of the rate.
\item Forgetting units in an interpretation.
\end{itemize}
\end{mistakes}

\begin{quickcheck}
1. The average rate of change of $g(x)=5-x^2$ on $[-1,2]$ is \fillin{$-1$}.\\
2. If $T(4)=21$ and $T(10)=9$ ($^\circ$C, hours), the average rate of change on $[4,10]$ is \fillin{$-2$ $^\circ$C per hour}.
\end{quickcheck}

%%DOC 1.3 | Rates of Change in Linear and Quadratic Functions
\goals{\item Recognize linear functions by their constant rate of change.
\item Recognize quadratic functions by their constant second differences, and find the quadratic from a table.
\item Connect the sign of the leading coefficient to concavity and to the vertex.}
\vocab{linear function, slope, quadratic function, first differences, second differences, vertex, axis of symmetry.}

\subsection*{1. Linear functions}
\begin{keyidea}
A linear function $f(x)=mx+b$ has the \fillin{same} average rate of change, equal to \fillin{$m$ (the slope)}, over every interval. In a table with equally spaced inputs, the outputs change by a \fillin{constant amount}.
\end{keyidea}

\subsection*{2. Quadratic functions}
\begin{keyidea}
For a quadratic $f(x)=ax^2+bx+c$, the average rates of change over consecutive equal-length intervals change at a \fillin{constant} rate. With equally spaced inputs, the \fillin{second differences} are constant. When the input step is 1, the second differences equal \fillin{$2a$}.
\end{keyidea}

\ex{Example 1 \textbullet\ Classify each table}
\begin{center}\begin{tabular}{c|cccc}$x$&0&1&2&3\\\hline A&4&7&10&13\\ B&1&2&5&10\\ C&1&2&4&8\end{tabular}\end{center}
\begin{workspace}[2.8cm]
A: differences $3,3,3$: linear. \quad B: differences $1,3,5$, second differences $2,2$: quadratic ($y=x^2+1$). \quad C: differences $1,2,4$, second differences $1,2$: neither linear nor quadratic (it doubles: exponential).
\end{workspace}

\ex{Example 2 \textbullet\ Finding the quadratic}
\begin{center}\begin{tabular}{c|ccccc}$x$&$0$&$1$&$2$&$3$&$4$\\\hline$y$&$3$&$1$&$-3$&$-9$&$-17$\end{tabular}\end{center}
Show that the data are quadratic and find the function.
\begin{workspace}[3.6cm]
First differences: $-2,-4,-6,-8$. Second differences: $-2,-2,-2$. Quadratic with $2a=-2$, so $a=-1$.\\
$c=y(0)=3$. Using $x=1$: $-1+b+3=1$, so $b=-1$. \quad $y=-x^2-x+3$.
\end{workspace}

\subsection*{3. Concavity and the vertex}
\begin{keyidea}
If $a>0$, the rate of change is \fillin{increasing}: the parabola is concave \fillin{up} and has a \fillin{minimum} at the vertex.\\
If $a<0$, it is concave \fillin{down} with a \fillin{maximum}. The vertex is at $x=$ \fillin{$-\frac{b}{2a}$}; there the rate of change changes sign.
\end{keyidea}

\ex{Example 3 \textbullet\ From a graph}
\begin{center}
\begin{tikzpicture}
\begin{axis}[gridaxis,xmin=-1.5,xmax=5.5,ymin=-1.5,ymax=5.5,width=6cm,height=5.5cm]
\addplot[apblue,very thick,domain=-0.5:4.5]{-(x-2)^2+5};
\node[circle,fill=apblue,inner sep=1.6pt] at (axis cs:2,5){};\node[circle,fill=apblue,inner sep=1.6pt] at (axis cs:0,1){};
\end{axis}
\end{tikzpicture}
\end{center}
Write an equation for this parabola and compute the average rates of change on $[0,1]$, $[1,2]$, $[2,3]$.
\begin{workspace}[3cm]
Vertex $(2,5)$: $f(x)=a(x-2)^2+5$. Through $(0,1)$: $4a+5=1$, $a=-1$. So $f(x)=-(x-2)^2+5$.\\
$f(0)=1$, $f(1)=4$, $f(2)=5$, $f(3)=4$: rates $3$, $1$, $-1$. They decrease by $2$ each step (concave down).
\end{workspace}

\begin{mistakes}
\begin{itemize}
\item Using differences when the inputs are not equally spaced.
\item Thinking the second difference equals $a$; it equals $2a$ when the step is 1.
\end{itemize}
\end{mistakes}

\begin{quickcheck}
1. Outputs $2,5,10,17$ at $x=0,1,2,3$ come from a \fillin{quadratic} function.\\
2. Second differences of $6$ (step 1) mean $a=$ \fillin{$3$}, and the graph is concave \fillin{up}.
\end{quickcheck}

%%DOC 1.4 | Polynomial Functions and Rates of Change
\goals{\item Identify the degree, leading term and leading coefficient of a polynomial in any form.
\item Locate local and global extrema from a graph, including on closed intervals.
\item Identify points of inflection and describe the rate of change between turning points.}
\vocab{polynomial, degree, leading term, leading coefficient, local (relative) extremum, global (absolute) extremum, point of inflection.}

\subsection*{1. Polynomial functions}
A polynomial of degree $n$ is $p(x)=a_nx^n+a_{n-1}x^{n-1}+\dots+a_0$ with $a_n$ \fillin{$\neq0$}. $a_nx^n$ is the \fillin{leading term} and $a_n$ the \fillin{leading coefficient}. The degree of a product is the \fillin{sum} of the degrees of the factors.

\ex{Example 1}
Give the degree and leading coefficient of $p(x)=-2x(x-1)^2(x+3)^3$.
\begin{workspace}[2.2cm]
Degree $1+2+3=6$. Leading term $-2x\cdot x^2\cdot x^3=-2x^6$: leading coefficient $-2$.
\end{workspace}

\subsection*{2. Extrema}
\begin{keyidea}
\begin{itemize}
\item A \textbf{local maximum} occurs where $p$ changes from \fillin{increasing to decreasing}.
\item A \textbf{local minimum} occurs where $p$ changes from \fillin{decreasing to increasing}.
\item The \textbf{global maximum} on an interval is the \fillin{greatest} output value on that interval; on a closed interval it can occur at an \fillin{endpoint}.
\item A polynomial of degree $n$ has at most \fillin{$n-1$} local extrema.
\end{itemize}
\end{keyidea}

\ex{Example 2 \textbullet\ Reading extrema}
\begin{center}
\begin{tikzpicture}
\begin{axis}[gridaxis,xmin=-3.5,xmax=3.5,ymin=-5.5,ymax=3.5,width=7cm,height=6cm]
\addplot[apblue,very thick,domain=-3.1:3.1,samples=150]{0.25*x^4-2*x^2};
\end{axis}
\end{tikzpicture}
\end{center}
The graph of $p(x)=\frac14x^4-2x^2$ is shown. Find all local extrema, the global extrema on $(-\infty,\infty)$, and the global maximum on $[0,3]$.
\begin{workspace}[3.2cm]
Local minima $(-2,-4)$ and $(2,-4)$; local maximum $(0,0)$.\\
Global minimum value $-4$; no global maximum (the ends go up).\\
On $[0,3]$: compare $p(0)=0$, $p(2)=-4$, $p(3)=\frac{81}{4}-18=2.25$. Global maximum $2.25$ at $x=3$ (an endpoint).
\end{workspace}

\subsection*{3. Points of inflection}
\begin{keyidea}
A point of inflection is where the graph changes \fillin{concavity}: the rate of change switches from increasing to decreasing, or the reverse. Between a local maximum and the next local minimum there is at least one point of inflection.
\end{keyidea}

\ex{Example 3}
$p(x)=x^3-6x^2+9x$ has a local maximum at $(1,4)$, a local minimum at $(3,0)$ and a point of inflection at $(2,2)$. Describe $p$ and its rate of change.
\begin{workspace}[3.2cm]
Increasing on $(-\infty,1)$, decreasing on $(1,3)$, increasing on $(3,\infty)$.\\
The rate of change decreases until $x=2$ (concave down on $(-\infty,2)$) and increases after (concave up on $(2,\infty)$). On $(1,2)$ the function decreases faster and faster; on $(2,3)$ it decreases more and more slowly.
\end{workspace}

\begin{mistakes}
\begin{itemize}
\item Giving the $x$-value when asked for the extreme \emph{value}, or the reverse.
\item Forgetting the endpoints when looking for global extrema on a closed interval.
\end{itemize}
\end{mistakes}

\begin{quickcheck}
1. A degree-6 polynomial has at most \fillin{5} turning points.\\
2. A cubic with no local extrema: \fillin{$p(x)=x^3$}.
\end{quickcheck}

%%DOC 1.5 | Polynomial Functions and Complex Zeros
\goals{\item Connect real zeros, linear factors and $x$-intercepts.
\item Use multiplicity to predict the graph near each zero.
\item Apply the Fundamental Theorem of Algebra and the conjugate pairs rule; write polynomials from their zeros.
\item Determine whether a function is even, odd, or neither.}
\vocab{zero, factor, multiplicity, complex number, imaginary unit $i$, conjugate, even function, odd function.}

\subsection*{1. Zeros, factors and intercepts}
\begin{keyidea}
For a polynomial $p$ and a real number $a$, these statements are equivalent: $a$ is a \fillin{zero} of $p$; $(x-a)$ is a \fillin{factor} of $p(x)$; $(a,0)$ is an \fillin{$x$-intercept} of the graph.
\end{keyidea}

\subsection*{2. Multiplicity}
If $(x-a)^k$ is a factor of $p(x)$ but $(x-a)^{k+1}$ is not, $a$ is a zero of multiplicity \fillin{$k$}.
\begin{keyidea}
Even multiplicity: the graph \fillin{touches the $x$-axis and turns back}.\\
Odd multiplicity: the graph \fillin{crosses the $x$-axis}; for multiplicity $3$ or more it also \fillin{flattens} as it crosses.
\end{keyidea}
\begin{center}
\begin{tikzpicture}
\begin{axis}[apaxis,xmin=-1.6,xmax=1.6,ymin=-1.6,ymax=1.6,width=4.4cm,height=4cm,xtick=\empty,ytick=\empty,title={\scriptsize multiplicity 1}]
\addplot[apblue,very thick,domain=-1.5:1.5]{x};\end{axis}\end{tikzpicture}
\begin{tikzpicture}
\begin{axis}[apaxis,xmin=-1.6,xmax=1.6,ymin=-1.6,ymax=1.6,width=4.4cm,height=4cm,xtick=\empty,ytick=\empty,title={\scriptsize multiplicity 2}]
\addplot[apblue,very thick,domain=-1.2:1.2]{x^2};\end{axis}\end{tikzpicture}
\begin{tikzpicture}
\begin{axis}[apaxis,xmin=-1.6,xmax=1.6,ymin=-1.6,ymax=1.6,width=4.4cm,height=4cm,xtick=\empty,ytick=\empty,title={\scriptsize multiplicity 3}]
\addplot[apblue,very thick,domain=-1.15:1.15]{x^3};\end{axis}\end{tikzpicture}
\end{center}

\subsection*{3. Complex zeros}
\begin{keyidea}
\textbf{Fundamental Theorem of Algebra}: a polynomial of degree $n$ has exactly \fillin{$n$} complex zeros, counted with multiplicity.\\
\textbf{Conjugate pairs}: if a polynomial with real coefficients has the zero $a+bi$, it also has the zero \fillin{$a-bi$}. Non-real zeros are \fillin{not} $x$-intercepts.
\end{keyidea}

\ex{Example 1 \textbullet\ Finding all zeros}
Find all zeros of $p(x)=x^4-5x^2-36$.
\begin{workspace}[2.6cm]
$x^4-5x^2-36=(x^2-9)(x^2+4)=(x-3)(x+3)(x^2+4)$. Zeros: $3$, $-3$ (real, $x$-intercepts) and $\pm2i$ (non-real).
\end{workspace}

\ex{Example 2 \textbullet\ Building a polynomial}
Write a polynomial of least degree with real coefficients, leading coefficient 1, and zeros $-1$ (multiplicity 2) and $2i$. Give its degree.
\begin{workspace}[2.6cm]
$-2i$ is also a zero. $p(x)=(x+1)^2(x-2i)(x+2i)=(x+1)^2(x^2+4)$, degree 4.
\end{workspace}

\ex{Example 3 \textbullet\ From a graph}
\begin{center}
\begin{tikzpicture}
\begin{axis}[gridaxis,xmin=-3.5,xmax=2.5,ymin=-3.5,ymax=5.5,width=5.5cm,height=5.5cm]
\addplot[apblue,very thick,domain=-2.6:2.1,samples=120]{(x+2)*(x-1)^2};
\end{axis}
\end{tikzpicture}
\end{center}
The graph of a cubic $q$ is shown. Write $q(x)$.
\begin{workspace}[2.6cm]
Crosses at $-2$ (odd multiplicity), touches at $1$ (even): $q(x)=a(x+2)(x-1)^2$. The $y$-intercept is $2$: $a\cdot2\cdot1=2$, $a=1$. So $q(x)=(x+2)(x-1)^2$.
\end{workspace}

\subsection*{4. Even and odd functions}
\begin{keyidea}
$f$ is \textbf{even} if $f(-x)=$ \fillin{$f(x)$} for every $x$; its graph is symmetric about the \fillin{$y$-axis}.\\
$f$ is \textbf{odd} if $f(-x)=$ \fillin{$-f(x)$}; its graph is symmetric about the \fillin{origin}.
\end{keyidea}

\ex{Example 4}
Classify $f(x)=x^4-3x^2$, $g(x)=x^3-5x$, $h(x)=x^3+1$.
\begin{workspace}[2.4cm]
$f(-x)=x^4-3x^2=f(x)$: even. \ $g(-x)=-x^3+5x=-g(x)$: odd. \ $h(-x)=-x^3+1$, equal to neither $h(x)$ nor $-h(x)$: neither.
\end{workspace}

\begin{mistakes}
\begin{itemize}
\item Listing non-real zeros as $x$-intercepts.
\item Writing the factor for the zero $-3$ as $(x-3)$.
\item Forgetting the conjugate when building a polynomial with real coefficients.
\end{itemize}
\end{mistakes}

\begin{quickcheck}
1. A degree-4 polynomial with real coefficients can have \fillin{0, 2 or 4} real zeros (with multiplicity).\\
2. $f(x)=x^5-2x^3+x$ is \fillin{odd}.
\end{quickcheck}

%%DOC 1.6 | Polynomial Functions and End Behavior
\goals{\item Describe end behavior with limit notation.
\item Determine end behavior from the degree and the sign of the leading coefficient, in standard or factored form.
\item Sketch a polynomial from its zeros, multiplicities and end behavior.}
\vocab{end behavior, limit, $\lim_{x\to\infty}$, $\lim_{x\to-\infty}$, increases without bound, leading term.}

\subsection*{1. Limit notation}
$\displaystyle\lim_{x\to\infty}p(x)=\infty$ means: as $x$ \fillin{increases without bound}, $p(x)$ \fillin{increases without bound}.\\
$\displaystyle\lim_{x\to-\infty}p(x)=-\infty$ means: as $x$ \fillin{decreases without bound}, $p(x)$ \fillin{decreases without bound}.

\subsection*{2. The leading term controls the ends}
\begin{keyidea}
\begin{center}\renewcommand{\arraystretch}{1.4}
\begin{tabular}{l|c|c}
& $a_n>0$ & $a_n<0$\\\hline
$n$ even & left \fillin{$\infty$}, right \fillin{$\infty$} & left \fillin{$-\infty$}, right \fillin{$-\infty$}\\\hline
$n$ odd & left \fillin{$-\infty$}, right \fillin{$\infty$} & left \fillin{$\infty$}, right \fillin{$-\infty$}
\end{tabular}\end{center}
\end{keyidea}

\ex{Example 1}
Describe the end behavior of $p(x)=3x^2-2x^5+7$.
\begin{workspace}[2.2cm]
Leading term $-2x^5$ (highest power, not the first term written): odd degree, negative. $\lim_{x\to-\infty}p(x)=\infty$, $\lim_{x\to\infty}p(x)=-\infty$.
\end{workspace}

\ex{Example 2 \textbullet\ Factored form}
Describe the end behavior of $q(x)=-x(x-2)^2(3-x)$.
\begin{workspace}[2.2cm]
Leading term: $-x\cdot x^2\cdot(-x)=x^4$. Even degree, positive: $\lim_{x\to\pm\infty}q(x)=\infty$.
\end{workspace}

\ex{Example 3 \textbullet\ Sketching}
Sketch $p(x)=(x+1)^2(x-2)$ using zeros, multiplicity, $y$-intercept and end behavior.
\begin{workspace}[5.2cm]
Zeros: $-1$ (multiplicity 2, touches), $2$ (crosses). $y$-intercept $p(0)=(1)(-2)=-2$. Leading term $x^3$: down on the left, up on the right.
\begin{center}
\begin{tikzpicture}
\begin{axis}[gridaxis,xmin=-2.5,xmax=3,ymin=-4.5,ymax=3.5,width=5cm,height=4.2cm]
\addplot[apblue,very thick,domain=-2.2:2.6,samples=100]{(x+1)^2*(x-2)};
\end{axis}
\end{tikzpicture}
\end{center}
\end{workspace}

\begin{mistakes}
\begin{itemize}
\item Using the first written term as the leading term.
\item In factored form, missing a negative sign hidden in a factor like $(3-x)$.
\end{itemize}
\end{mistakes}

\begin{quickcheck}
1. $\lim_{x\to-\infty}(4-x^3)=$ \fillin{$\infty$}.\\
2. Both ends go down when the degree is \fillin{even} and the leading coefficient is \fillin{negative}.
\end{quickcheck}

%%DOC 1.7 | Rational Functions and End Behavior
\goals{\item Determine the end behavior of a rational function by comparing the degrees of numerator and denominator.
\item Find horizontal and slant asymptotes, and write end behavior with limits.}
\vocab{rational function, horizontal asymptote, slant asymptote, quotient of leading terms.}

\subsection*{1. Quotient of leading terms}
For $r(x)=\dfrac{p(x)}{q(x)}$, as $x\to\pm\infty$ the function behaves like \fillin{$\dfrac{\text{leading term of }p}{\text{leading term of }q}$}.

\begin{keyidea}
\begin{tabular}{@{}ll@{}}
$\deg p<\deg q$: & horizontal asymptote \fillin{$y=0$}\\[2pt]
$\deg p=\deg q$: & horizontal asymptote $y=$ \fillin{ratio of the leading coefficients}\\[2pt]
$\deg p=\deg q+1$: & \fillin{slant} asymptote, found by polynomial division\\[2pt]
$\deg p>\deg q+1$: & behaves like a \fillin{polynomial} of degree $\deg p-\deg q$
\end{tabular}
\end{keyidea}

\ex{Example 1 \textbullet\ Limits at infinity}
Find $\lim_{x\to\infty}$ for $\dfrac{3x+1}{x^2-4}$, \ $\dfrac{8x^2-x}{2x^2+5}$, \ $\dfrac{x^3}{x^2+1}$.
\begin{workspace}[2.6cm]
$0$ (degree of the denominator is larger). \quad $\frac82=4$. \quad Behaves like $x$: $\infty$ (and as $x\to-\infty$, $-\infty$).
\end{workspace}

\ex{Example 2 \textbullet\ Slant asymptote}
Find the slant asymptote of $s(x)=\dfrac{x^2+x-1}{x-1}$.
\begin{workspace}[5cm]
$x^2+x-1=(x-1)(x+2)+1$, so $s(x)=x+2+\dfrac{1}{x-1}$. As $x\to\pm\infty$, $\frac{1}{x-1}\to0$: slant asymptote $y=x+2$.
\begin{center}
\begin{tikzpicture}
\begin{axis}[gridaxis,xmin=-4.5,xmax=5.5,ymin=-4.5,ymax=8.5,ytick distance=2,width=5.5cm,height=4.4cm]
\addplot[apblue,very thick,domain=-4.5:0.8,samples=150]{x+2+1/(x-1)};
\addplot[apblue,very thick,domain=1.15:5.5,samples=150]{x+2+1/(x-1)};
\addplot[apgold!80!black,dashed,thick,domain=-4.5:5.5]{x+2};
\draw[apred,dashed](axis cs:1,-10)--(axis cs:1,10);
\end{axis}
\end{tikzpicture}
\end{center}
\end{workspace}

\ex{Example 3 \textbullet\ From a graph}
\begin{center}
\begin{tikzpicture}
\begin{axis}[gridaxis,xmin=-6.5,xmax=6.5,ymin=-3.5,ymax=1.5,width=7cm,height=4cm]
\addplot[apblue,very thick,domain=-6.5:6.5,samples=150]{-2*x^2/(x^2+1)};
\draw[apgold!80!black,dashed,thick](axis cs:-7,-2)--(axis cs:7,-2);
\end{axis}
\end{tikzpicture}
\end{center}
The graph of $r(x)=\frac{-2x^2}{x^2+1}$ is shown. Write the end behavior with limits.
\begin{workspace}[1.8cm]
$\lim_{x\to\infty}r(x)=-2$ and $\lim_{x\to-\infty}r(x)=-2$ (horizontal asymptote $y=-2$).
\end{workspace}

\begin{mistakes}
\begin{itemize}
\item Reading the asymptote from the constant terms.
\item Believing a graph can never cross a horizontal asymptote: it can, for moderate $x$.
\end{itemize}
\end{mistakes}

\begin{quickcheck}
1. Horizontal asymptote of $\dfrac{5-2x^3}{x^3+x}$: \fillin{$y=-2$}.\\
2. Slant asymptote of $\dfrac{2x^2+3x}{x+1}$: \fillin{$y=2x+1$}.
\end{quickcheck}

%%DOC 1.8 | Rational Functions and Zeros
\goals{\item Find the real zeros of a rational function, checking the domain.
\item Solve rational inequalities with a sign chart.}
\vocab{zero, domain, excluded value, sign chart, test value, interval notation.}

\subsection*{1. Zeros}
\begin{keyidea}
The real zeros of $r=\dfrac pq$ are the real zeros of \fillin{$p$} that are \fillin{in the domain of $r$} (not zeros of $q$). Always \fillin{factor} first.
\end{keyidea}

\ex{Example 1}
Find the zeros of $r(x)=\dfrac{x^2-3x}{x^2-9}$ and of $s(x)=\dfrac{x^2-4}{x+1}$.
\begin{workspace}[2.6cm]
$r(x)=\frac{x(x-3)}{(x-3)(x+3)}$: numerator zeros $0,3$, but $x=3$ is excluded. Zero: $x=0$ only.\\
$s$: zeros $x=2$ and $x=-2$ (neither is excluded).
\end{workspace}

\subsection*{2. Sign charts}
The sign of a rational function can change only at its \fillin{zeros} and at \fillin{values excluded from its domain}. Mark these on a number line, test \fillin{one value in each interval}, and read off the sign.

\ex{Example 2 \textbullet\ Solve $r(x)\le0$}
Solve $\dfrac{(x+1)(x-4)}{x-2}\le0$.
\begin{workspace}[4.6cm]
Key values: $-1$, $2$, $4$.
\begin{center}
\begin{tikzpicture}[x=1cm]
\draw[->,apgrey](-2.5,0)--(5.5,0);
\foreach \p in {-1,2,4}{\draw (\p,0.12)--(\p,-0.12) node[below]{\small$\p$};}
\node[apred] at (-1.8,0.35){$-$};\node[apgreen] at (0.5,0.35){$+$};\node[apred] at (3,0.35){$-$};\node[apgreen] at (4.8,0.35){$+$};
\node[circle,fill=apnavy,inner sep=1.5pt] at (-1,0){};\node[circle,draw=apnavy,fill=white,inner sep=1.5pt] at (2,0){};\node[circle,fill=apnavy,inner sep=1.5pt] at (4,0){};
\end{tikzpicture}
\end{center}
Tests: $x=-2$: negative; $x=0$: positive; $x=3$: negative; $x=5$: positive. Include zeros, exclude $x=2$: $(-\infty,-1]\cup(2,4]$.
\end{workspace}

\ex{Example 3 \textbullet\ From a graph}
\begin{center}
\begin{tikzpicture}
\begin{axis}[gridaxis,xmin=-7.5,xmax=7.5,ymin=-8.5,ymax=3.5,width=7.5cm,height=5cm]
\addplot[apblue,very thick,domain=-7.5:-0.6,samples=200]{(x+3)*(x-2)/x^2};
\addplot[apblue,very thick,domain=0.55:7.5,samples=200]{(x+3)*(x-2)/x^2};
\draw[apred,dashed,thick](axis cs:0,-10)--(axis cs:0,10);
\end{axis}
\end{tikzpicture}
\end{center}
The graph of $r(x)=\frac{(x+3)(x-2)}{x^2}$ is shown. Solve $r(x)>0$.
\begin{workspace}[1.8cm]
The graph is above the axis for $x<-3$ and $x>2$: $(-\infty,-3)\cup(2,\infty)$.
\end{workspace}

\begin{mistakes}
\begin{itemize}
\item Including excluded values in the solution of $\le$ or $\ge$.
\item Multiplying both sides of an inequality by an expression whose sign is unknown.
\end{itemize}
\end{mistakes}

\begin{quickcheck}
1. Zeros of $\dfrac{x^2-1}{x+1}$: \fillin{$x=1$ only}.\qquad 2. $\dfrac{x}{x-3}>0$ on \fillin{$(-\infty,0)\cup(3,\infty)$}.
\end{quickcheck}

%%DOC 1.9-1.10 | Rational Functions and Asymptotes
\goals{\item Locate vertical asymptotes and holes of a rational function from its factored form.
\item Describe behavior near a vertical asymptote with one-sided limits, determining signs.
\item Find the coordinates of a hole.}
\vocab{vertical asymptote, hole, one-sided limit, $\lim_{x\to a^-}$, $\lim_{x\to a^+}$, multiplicity, simplified form.}

\subsection*{1. Vertical asymptotes or holes?}
\begin{keyidea}
Factor numerator and denominator completely. For a zero $a$ of the denominator:
\begin{itemize}
\item if $(x-a)$ appears \fillin{more} times in the denominator than in the numerator, $x=a$ is a \fillin{vertical asymptote};
\item if it appears \fillin{at least as many} times in the numerator, there is a \fillin{hole} at $x=a$.
\end{itemize}
\end{keyidea}

\subsection*{2. Behavior near a vertical asymptote}
Near $x=a$, $|r(x)|$ grows \fillin{without bound}. To decide $+\infty$ or $-\infty$, check the \fillin{sign} of each factor just left and just right of $a$. If $(x-a)$ remains an \fillin{odd} number of times in the denominator, the sign flips across the asymptote; an \fillin{even} number of times, it stays the same.

\ex{Example 1}
Describe $r(x)=\dfrac{x+1}{(x-2)^2}$ and $s(x)=\dfrac{x-1}{(x+2)(x-3)}$ near their vertical asymptotes.
\begin{workspace}[4cm]
$r$: numerator near $3>0$, $(x-2)^2>0$ on both sides: $\lim_{x\to2^-}r=\lim_{x\to2^+}r=+\infty$.\\
$s$ at $-2$: numerator near $-3$, $x-3$ near $-5$. For $x<-2$ the denominator is positive, so $s<0$: $-\infty$; for $x>-2$: $+\infty$.\\
$s$ at $3$: numerator near $2$, $x+2$ near $5$. Left: $-\infty$; right: $+\infty$.
\end{workspace}

\subsection*{3. Holes}
\begin{keyidea}
To find the hole at $x=a$: cancel the common factors, then substitute $a$ into the \fillin{simplified} expression. The hole is at $\big(a,\ \text{that value}\big)$, and $\lim_{x\to a}r(x)$ equals \fillin{that value}.
\end{keyidea}

\ex{Example 2 \textbullet\ Full analysis}
Analyze $r(x)=\dfrac{x^2-x-6}{x^2-9}$: hole, vertical asymptote, zero, end behavior. Then sketch.
\begin{workspace}[6cm]
$r(x)=\dfrac{(x-3)(x+2)}{(x-3)(x+3)}$. Hole at $x=3$: $\frac{5}{6}$, so $\left(3,\frac56\right)$. Vertical asymptote $x=-3$: $\lim_{x\to-3^-}r=+\infty$, $\lim_{x\to-3^+}r=-\infty$. Zero $x=-2$. Horizontal asymptote $y=1$.
\begin{center}
\begin{tikzpicture}
\begin{axis}[gridaxis,xmin=-8,xmax=7,ymin=-5,ymax=6,xtick distance=2,width=6cm,height=4.6cm]
\addplot[apblue,very thick,domain=-8:-3.18,samples=200]{(x+2)/(x+3)};
\addplot[apblue,very thick,domain=-2.83:7,samples=200]{(x+2)/(x+3)};
\draw[apred,dashed,thick] (axis cs:-3,-6)--(axis cs:-3,7);\draw[apgold!80!black,dashed,thick] (axis cs:-8,1)--(axis cs:7,1);
\node[circle,draw=apblue,fill=white,inner sep=1.6pt,thick] at (axis cs:3,5/6){};
\end{axis}
\end{tikzpicture}
\end{center}
\end{workspace}

\ex{Example 3 \textbullet\ A hole on the axis}
$r(x)=\dfrac{(x-1)^2}{x-1}$. Is $x=1$ a zero?
\begin{workspace}[2cm]
No. $r(x)=x-1$ for $x\neq1$, but $r(1)$ is undefined. The graph has a hole at $(1,0)$, not a zero.
\end{workspace}

\begin{aptip}
Justify each feature with the factored form \emph{and} a limit: ``$\lim_{x\to-3^+}r(x)=-\infty$, so $x=-3$ is a vertical asymptote''; ``$\lim_{x\to3}r(x)=\frac56$ and $r(3)$ is undefined, so there is a hole at $\left(3,\frac56\right)$.''
\end{aptip}

\begin{mistakes}
\begin{itemize}
\item Calling every zero of the denominator a vertical asymptote without checking for cancellation.
\item Finding the hole's height in the original expression ($\frac00$).
\end{itemize}
\end{mistakes}

\begin{quickcheck}
1. $\dfrac{x-3}{x^2-9}$ has a vertical asymptote at \fillin{$x=-3$} and a hole at \fillin{$\left(3,\frac16\right)$}.\\
2. $\lim_{x\to0^-}\dfrac{2}{x^3}=$ \fillin{$-\infty$}.
\end{quickcheck}

%%DOC 1.11 | Rational Functions and Concept of Limit
\goals{\item Use limit notation to describe the behavior of rational functions near holes, near vertical asymptotes, and at the ends.
\item Rewrite rational expressions (factoring, polynomial division) to evaluate limits and find asymptotes.
\item Use the remainder theorem and the binomial theorem to work with equivalent forms of polynomials.}
\vocab{limit, one-sided limit, equivalent expression, polynomial division, remainder theorem, binomial theorem, Pascal's triangle.}

\subsection*{1. Three kinds of limit statements}
\begin{keyidea}
\begin{tabular}{@{}lll@{}}
Hole at $x=a$: & $\displaystyle\lim_{x\to a}r(x)=$ \fillin{a real number $L$} & ($r(a)$ undefined)\\[6pt]
Vertical asymptote $x=a$: & $\displaystyle\lim_{x\to a^\pm}r(x)=$ \fillin{$\pm\infty$} & \\[6pt]
Horizontal asymptote $y=b$: & $\displaystyle\lim_{x\to\pm\infty}r(x)=$ \fillin{$b$} &
\end{tabular}
\end{keyidea}

\ex{Example 1 \textbullet\ Describe everything with limits}
$r(x)=\dfrac{2x^2+2x-4}{x^2-1}$.
\begin{workspace}[3.6cm]
$r(x)=\dfrac{2(x+2)(x-1)}{(x-1)(x+1)}$. Hole: $\lim_{x\to1}r(x)=\frac{2(3)}{2}=3$. Asymptote: $\lim_{x\to-1^-}r(x)=-\infty$, $\lim_{x\to-1^+}r(x)=+\infty$ (numerator of $\frac{2(x+2)}{x+1}$ near $2$). Ends: $\lim_{x\to\pm\infty}r(x)=2$.
\end{workspace}

\subsection*{2. Equivalent forms reveal limits}
\begin{keyidea}
Dividing $p(x)$ by $(x-a)$ gives $p(x)=(x-a)q(x)+R$, and the \textbf{remainder theorem} says $R=$ \fillin{$p(a)$}. So $(x-a)$ is a factor exactly when \fillin{$p(a)=0$}.\\
Writing $\dfrac{p(x)}{d(x)}=q(x)+\dfrac{R(x)}{d(x)}$ shows that the graph approaches $y=q(x)$ as $x\to\pm\infty$ when $\deg R<\deg d$.
\end{keyidea}

\ex{Example 2 \textbullet\ Synthetic division}
Divide $2x^3-3x^2-11x+6$ by $x-3$ and factor completely.
\begin{workspace}[3.4cm]
\begin{tabular}{r|rrrr}$3$&$2$&$-3$&$-11$&$6$\\&&$6$&$9$&$-6$\\\hline&$2$&$3$&$-2$&$0$\end{tabular}\qquad
Quotient $2x^2+3x-2=(2x-1)(x+2)$, remainder $0$.\\
$2x^3-3x^2-11x+6=(x-3)(2x-1)(x+2)$; zeros $3,\ \frac12,\ -2$.
\end{workspace}

\ex{Example 3 \textbullet\ A limit and an asymptote from division}
Let $r(x)=\dfrac{x^3+2x^2-5x+1}{x-1}$. Find $\lim_{x\to1^+}r(x)$ and describe the end behavior.
\begin{workspace}[3cm]
Synthetic division by $1$: $1,3,-2$, remainder $-1$: $r(x)=x^2+3x-2-\dfrac{1}{x-1}$.\\
As $x\to1^+$, $-\frac{1}{x-1}\to-\infty$: $\lim_{x\to1^+}r(x)=-\infty$. As $x\to\pm\infty$, $r$ behaves like $x^2+3x-2$: $\lim_{x\to\pm\infty}r(x)=\infty$.
\end{workspace}

\subsection*{3. The binomial theorem}
The coefficients of $(a+b)^n$ are row $n$ of \fillin{Pascal's triangle}: row 4 is \fillin{$1,4,6,4,1$}.
\ex{Example 4}
Expand $(x-2)^4$ and $(2x+1)^3$.
\begin{workspace}[2.4cm]
$(x-2)^4=x^4-8x^3+24x^2-32x+16$.\qquad $(2x+1)^3=8x^3+12x^2+6x+1$.
\end{workspace}

\begin{mistakes}
\begin{itemize}
\item Writing $\lim_{x\to a}r(x)=\infty$ when the one-sided limits have opposite signs.
\item Leaving out zero coefficients in synthetic division ($x^3+5$ is $1,0,0,5$).
\end{itemize}
\end{mistakes}

\begin{quickcheck}
1. The remainder when $x^4-3x+1$ is divided by $x+1$ is \fillin{$5$}.\\
2. $\lim_{x\to2}\dfrac{x^2-4}{x-2}=$ \fillin{$4$}. \quad 3. The coefficient of $x^2$ in $(x+3)^3$ is \fillin{$9$}.
\end{quickcheck}

%%DOC 1.12 | Transformations of Functions
\goals{\item Describe translations, dilations and reflections, and their order.
\item Write $g$ in terms of $f$ from a description, a table or a graph.
\item Find the domain and range of a transformed function and the image of a point.}
\vocab{translation, vertical dilation, horizontal dilation, reflection, image of a point, domain, range.}

\subsection*{1. The four parameters}
\begin{keyidea}
For $g(x)=a\,f\big(b(x-h)\big)+k$:
\begin{itemize}
\item $k$: vertical translation, \fillin{up} if $k>0$. \quad $h$: horizontal translation, \fillin{right} if $h>0$.
\item $a$: vertical dilation by a factor of \fillin{$|a|$}; if $a<0$, also a reflection over the \fillin{$x$-axis}.
\item $b$: horizontal dilation by a factor of \fillin{$\frac1{|b|}$}; if $b<0$, also a reflection over the \fillin{$y$-axis}.
\end{itemize}
A point $(x,y)$ on $f$ moves to \fillin{$\left(\frac{x}{b}+h,\ ay+k\right)$} on $g$. Horizontal changes affect the \fillin{domain}; vertical changes affect the \fillin{range}.
\end{keyidea}

\ex{Example 1 \textbullet\ Describe the transformations}
Describe how to obtain $g(x)=-2(x+3)^2+5$ from $f(x)=x^2$.
\begin{workspace}[2.2cm]
Shift left 3, vertical dilation by a factor of 2, reflection over the $x$-axis, shift up 5. The vertex moves from $(0,0)$ to $(-3,5)$.
\end{workspace}

\ex{Example 2 \textbullet\ From a graph}
\begin{center}
\begin{tikzpicture}
\begin{axis}[gridaxis,xmin=-1.5,xmax=6.5,ymin=-0.5,ymax=6.5,width=6cm,height=5cm]
\addplot[apblue,very thick] coordinates {(0,0)(1,2)(2,0)};
\addplot[apred,very thick,dashed] coordinates {(3,1)(4,5)(5,1)};
\node[apblue,font=\small] at (axis cs:0.6,2.4){$f$};\node[apred,font=\small] at (axis cs:4.6,5.4){$g$};
\end{axis}
\end{tikzpicture}
\end{center}
Write $g$ in terms of $f$.
\begin{workspace}[2.4cm]
Same width; height above the base doubles ($2\to4$); moved right 3 and up 1: $g(x)=2f(x-3)+1$. Check: $2f(1)+1=5$ at $x=4$.
\end{workspace}

\ex{Example 3 \textbullet\ Domain, range and points}
$f$ has domain $[-4,2]$ and range $[-1,3]$, and $(2,3)$ is on its graph. For $g(x)=-f(2x)+1$, find the domain, the range and the image of $(2,3)$.
\begin{workspace}[2.6cm]
Domain: $2x\in[-4,2]$, so $[-2,1]$. Range: $-f\in[-3,1]$, then $+1$: $[-2,2]$.\\
Image of $(2,3)$: $x=\frac22=1$, $y=-3+1=-2$: $(1,-2)$.
\end{workspace}

\ex{Example 4 \textbullet\ From a table}
\begin{center}\begin{tabular}{c|ccccc}$x$&$-2$&$-1$&$0$&$1$&$2$\\\hline$f(x)$&$5$&$2$&$1$&$2$&$5$\\$g(x)$&$2$&$1$&$2$&$5$&$10$\end{tabular}\end{center}
Write $g$ in terms of $f$.
\begin{workspace}[2cm]
The values of $g$ are those of $f$ one input earlier: $g(-2)=f(-1)$, $g(-1)=f(0)$, $g(0)=f(1)$, $g(1)=f(2)$. So $g(x)=f(x+1)$: the graph of $f$ shifted 1 unit left.
\end{workspace}

\begin{mistakes}
\begin{itemize}
\item Shifting the wrong way: $f(x+3)$ moves the graph \emph{left}.
\item Thinking $f(2x)$ stretches horizontally; it compresses by a factor of $\frac12$.
\end{itemize}
\end{mistakes}

\begin{quickcheck}
1. $(2,5)$ is on $f$. Then \fillin{$(1,2)$} is on $y=f(x+1)-3$.\\
2. $y=-\frac12f(x)$ is a vertical dilation by \fillin{$\frac12$} and a reflection over the \fillin{$x$-axis}.
\end{quickcheck}

%%DOC 1.13-1.14 | Function Model Construction
\goals{\item Choose an appropriate function type for data or a context, and state the assumptions.
\item Build models from geometric or verbal descriptions, and from data with regression.
\item Use models to predict and interpret, with units and a sensible domain.}
\vocab{model, linear, quadratic, cubic, rational, piecewise, regression, assumption, restricted domain, extrapolation.}

\subsection*{1. Choosing a model}
\begin{keyidea}
\begin{tabular}{@{}ll@{}}
Roughly constant rate of change & \fillin{linear}\\
Constant second differences; one maximum or minimum & \fillin{quadratic}\\
Two turning points or a change in concavity & \fillin{cubic}\\
Approaches a limiting value; ``per unit'' quantities & \fillin{rational}\\
Different rules on different intervals & \fillin{piecewise}
\end{tabular}
\end{keyidea}

\ex{Example 1 \textbullet\ Which model?}
\begin{center}
\begin{tikzpicture}
\begin{axis}[gridaxis,xmin=-0.5,xmax=6.5,ymin=-0.5,ymax=6.5,width=5.5cm,height=4.6cm,xlabel={hours},ylabel={}]
\addplot[only marks,mark=*,mark size=2pt,apblue] coordinates {(0,0.5)(1,2.4)(2,3.6)(3,4.3)(4,4.7)(5,4.9)(6,5.0)};
\end{axis}
\end{tikzpicture}
\end{center}
The concentration of a medicine rises, then levels off near 5. Which model fits?
\begin{workspace}[2cm]
The data approach a limiting value: a rational model with horizontal asymptote $y=5$ fits. A polynomial would eventually grow without bound.
\end{workspace}

\subsection*{2. Assumptions and domain}
Every model assumes the pattern \fillin{continues}. Restrict the domain to inputs that make sense in context (for example, time $\ge0$, lengths \fillin{positive}). Predictions far outside the data (\fillin{extrapolation}) are less reliable.

\subsection*{3. Building a model from a description}
\ex{Example 2 \textbullet\ Optimization}
A farmer has 60 m of fence for a rectangular pen along a wall (no fence on the wall side). Model the area and find the largest area.
\begin{workspace}[4.6cm]
Let $x$ = length of each side perpendicular to the wall. The parallel side is $60-2x$. $A(x)=x(60-2x)$, domain $0<x<30$.\\
$A(x)=-2(x-15)^2+450$: maximum area $450$ m$^2$ when $x=15$ m (and the parallel side is $30$ m).
\end{workspace}

\ex{Example 3 \textbullet\ Piecewise}
A phone plan costs \$20 per month including 100 texts, plus \$0.10 per extra text. Write $C(n)$ and find $C(250)$.
\begin{workspace}[2.6cm]
$C(n)=20$ for $0\le n\le100$; $C(n)=20+0.10(n-100)$ for $n>100$. $C(250)=35$: 250 texts cost \$35.
\end{workspace}

\subsection*{4. Building a model from data}
\ex{Example 4 \textbullet\ Regression (calculator)}
A ball's height (m) at $t=0,1,2,3,4$ s is $1.2,\ 6.1,\ 9.0,\ 9.9,\ 8.6$. Find a quadratic model and predict the height at $t=5$.
\begin{workspace}[2.6cm]
Quadratic regression: $h(t)\approx-1.029t^2+5.974t+1.183$. $h(5)\approx5.3$ m. Maximum about $9.9$ m at $t\approx2.9$ s.
\end{workspace}

\begin{aptip}
Keep regression coefficients stored in the calculator and round only the final answer (three decimal places). Always state units and the meaning in context.
\end{aptip}

\begin{mistakes}
\begin{itemize}
\item Leaving two variables in a model: use the constraint to eliminate one.
\item Using a model for inputs outside its domain (negative times, negative lengths).
\end{itemize}
\end{mistakes}

\begin{quickcheck}
1. ``The value drops by \$1200 each year'' suggests a \fillin{linear} model.\\
2. A box with a square base of side $x$ and volume 32 has height \fillin{$\frac{32}{x^2}$}.
\end{quickcheck}
